\documentclass[11pt,a4paper]{article}
\usepackage[margin=1in]{geometry}
\usepackage{amsmath,amssymb,amsthm}
\usepackage{mathrsfs}
\usepackage{enumitem}
\usepackage{hyperref}
\hypersetup{colorlinks=true, linkcolor=blue, citecolor=blue, urlcolor=blue}

% Theorem environments
\theoremstyle{plain}
\newtheorem{theorem}{Theorem}[section]
\newtheorem{corollary}{Corollary}[theorem]
\newtheorem{lemma}{Lemma}[section]
\theoremstyle{definition}
\newtheorem{definition}{Definition}[section]
\newtheorem{example}{Example}[section]
\newtheorem{application}{Application}[section]

\title{Four Fundamental Theorems of Functional Analysis}
% \author{Algorithm Research Institute}
\date{\today}

\begin{document}
\maketitle




\textbf{Intuition and Terminology for Beginners:}
Before diving into the theorem, let's break down the jargon:
\begin{itemize}
    \item \textbf{Evolution Equation:} A PDE that describes how a system changes over time \(t\) (e.g., heat diffusing or waves propagating). We know the initial state \(u_0\) and want to find the future state \(u(t)\).
    \item \textbf{Why ``Semigroup"?} Let \(T(t)\) be the ``time machine" operator that moves the system forward by \(t\) seconds. If we wait \(s\) seconds and then \(t\) seconds, it is the same as waiting \(t+s\) seconds, so \(T(t)T(s) = T(t+s)\). We call it a ``semi"-group because for equations like heat conduction, time is irreversible (\(t \ge 0\)), so we cannot go backward in time.
    \item \textbf{Why ``\(C_0\)"?} The ``\(C\)" stands for Continuous, and ``\(0\)" stands for time \(t=0\). Physically, it means the system cannot teleport or mutate instantly. If almost no time has passed (\(t \to 0\)), the state \(T(t)u_0\) must be exactly the initial state \(u_0\).
\end{itemize}
\begin{application}[Uniform Boundedness of $C_0$-Semigroups in Evolution PDEs]
    Let \(X\) be a Banach space. Consider the initial value problem for an evolution equation (such as the heat or wave equation)
    \[
        \frac{du}{dt} = Au, \quad u(0) = u_0 \in X,
    \]
    where \(A\) is a linear operator. Suppose the solution is given by a strongly continuous semigroup (\(C_0\)-semigroup) of bounded linear operators \(\{T(t)\}_{t\ge 0}\), such that \(u(t) = T(t)u_0\).

    By the definition of strong continuity, for every \(u_0 \in X\), the orbit \(t \mapsto T(t)u_0\) is continuous at \(t=0\), i.e.,
    \[
        \lim_{t \to 0^+} \|T(t)u_0 - u_0\|_X = 0.
    \]
    Then, the operator family \(\{T(t)\}\) is locally uniformly bounded. Specifically, there exists a constant \(M > 0\) such that
    \[
        \sup_{t \in [0,1]} \|T(t)\| \le M.
    \]
\end{application}
\textbf{Intuition and Terminology for Beginners:}
Why do we care about these two specific problems? They perfectly illustrate the power---and the strict requirements---of the Uniform Boundedness Principle (Banach-Steinhaus Theorem).
\begin{itemize}
    \item \textbf{Problem 9 (The Counterexample):} The theorem says ``pointwise bounded implies uniformly bounded." But this ONLY works in a complete space (Banach space). Problem 9 builds a space of polynomials that is not complete. We define an operator that truncates and sums the coefficients. It is pointwise bounded for every individual polynomial, but the overall operator norm explodes to infinity. The theorem fails spectacularly because the safety net (completeness) is missing!
    \item \textbf{Problem 12 (The Resonance Phenomenon):} This is one of the most famous applications of the theorem. We try to approximate a continuous function using Lagrange interpolation. We might intuitively think that using more nodes ($n \to \infty$) gives a better approximation. However, the ``magnification power" (Lebesgue constants) of the interpolation operator grows logarithmically. Because the space of continuous functions is strictly complete, the Resonance Theorem guarantees that there is at least one perfectly continuous function whose interpolation polynomials oscillate wildly and diverge to infinity!
\end{itemize}

\vspace{1.5em}
\hrule
\vspace{1.5em}

\begin{example}[Failure of Banach-Steinhaus in Incomplete Spaces]
    Let $P$ denote the space of all real polynomials[cite: 24]. Given a polynomial
    \[
        x(t) = \sum_{k=0}^{m}\alpha_{k}t^{k},
    \]
    we define the norm on $P$ as:
    \[
        \|x\| = \max_{0\le k\le m}|\alpha_{k}| \quad \text{[cite: 25]}.
    \]
    Define a sequence of linear functionals $T_n : P \to \mathbb{R}$ by truncating the series:
    \[
        T_{n}x = \sum_{k=0}^{\min\{m,n\}}\alpha_{k} \quad \text{[cite: 26]}.
    \]
    Then, the sequence is pointwise bounded, meaning
    \[
        \sup_{n\ge 0}|T_{n}x| < \infty \quad \forall x \in P \quad \text{[cite: 27]}.
    \]
    However, the sequence of operator norms is unbounded, i.e.,
    \[
        \sup_{n\ge 0}\|T_{n}\| = \infty \quad \text{[cite: 28]}.
    \]
    \textit{Remark:} This explicitly shows that the Banach-Steinhaus theorem fails when the space $X$ is not complete[cite: 29].
\end{example}

\begin{proof}
    \textbf{Step 1: Pointwise Boundedness} \medskip

    Fix any polynomial $x \in P$. It has a finite, fixed degree $m$. For all \(n \ge m\), the index bound becomes $\min\{m,n\} = m$. Thus, for all sufficiently large $n$, the functional evaluates to a constant:
    \[
        T_n x = \sum_{k=0}^{m}\alpha_{k}.
    \]
    Since the sequence $\{T_n x\}$ becomes constant for large $n$, it is strictly bounded. Thus, we have:
    \[
        \sup_{n\ge 0}|T_{n}x| < \infty \quad \text{[cite: 27]}.
    \]

    \textbf{Step 2: Unbounded Operator Norms} \medskip

    For any $x \in P$, we can estimate the upper bound of the functional:
    \[
        \begin{aligned}
            |T_n x| & \le \sum_{k=0}^{n}|\alpha_k|            \\
                    & \le (n+1) \max_{0\le k\le n} |\alpha_k| \\
                    & \le (n+1)\|x\|.
        \end{aligned}
    \]
    This implies that the operator norm is bounded above by $n+1$, i.e., $\|T_n\| \le n+1$.

    To show that this supremum is actually attained, consider the specific polynomial:
    \[
        x_n(t) = 1 + t + t^2 + \dots + t^n.
    \]
    Clearly, $x_n \in P$ and its norm is $\|x_n\| = \max\{1, 1, \dots, 1\} = 1$.
    Applying our operator $T_n$ to this polynomial yields:
    \[
        T_n x_n = \sum_{k=0}^{n} 1 = n+1.
    \]
    Therefore, the operator norm is exactly $\|T_n\| = n+1$.
    As $n \to \infty$, the norm $\|T_n\| \to \infty$, which rigorously proves:
    \[
        \sup_{n\ge 0}\|T_{n}\| = \infty \quad \text{[cite: 28]}.
    \]
\end{proof}

\vspace{1.5em}
\hrule
\vspace{1.5em}

\begin{application}[Divergence of Lagrange Interpolation]
    For each integer $n\ge 0$, let $0\le t_{0} < t_{1} < \dots < t_{n}\le 1$ be any set of $(n+1)$ distinct nodes[cite: 64]. The Lagrange interpolating operator $L_{n}:C[0,1]\rightarrow P[0,1]$ is defined by:
    \[
        L_{n}x(t_{k})=x(t_{k}), \quad 0\le k\le n, \quad \text{and} \quad \deg L_{n}x \le n \quad \text{[cite: 65, 66, 67]}.
    \]
    This interpolating polynomial can be explicitly written in the Lagrange basis form:
    \[
        L_{n}x(t) = \sum_{k=0}^{n}x(t_{k})p_{k}(t) \quad \text{[cite: 69]},
    \]
    where the basis polynomials are given by:
    \[
        p_{k}(t) = \prod_{j\ne k}^{n}\frac{t-t_{j}}{t_{k}-t_{j}} \quad \text{[cite: 70]}.
    \]
    We will establish two main results:
    \begin{enumerate}
        \item The operator norm (the Lebesgue constant) is exactly $\|L_{n}\| = \sup_{0\le t\le 1}\sum_{k=0}^{n}|p_{k}(t)|$[cite: 73].
        \item Given the known analytical result that $\|L_{n}\| \ge C \log n$ for some constant $C>0$[cite: 84, 85], there exists a continuous function $x\in C[0,1]$ such that $\sup_{n\ge 0}\|L_{n}x\| = \infty$[cite: 86, 87].
    \end{enumerate}
\end{application}

\begin{proof}
    \textbf{Part (i): Calculating the Operator Norm} \medskip

    Let $x \in C[0,1]$ with uniform norm $\|x\|_\infty \le 1$. We estimate the interpolation:
    \[
        \begin{aligned}
            |L_n x(t)| & \le \sum_{k=0}^{n} |x(t_k)| |p_k(t)|            \\
                       & \le \|x\|_\infty \sum_{k=0}^{n} |p_k(t)|        \\
                       & \le \sup_{0\le t\le 1} \sum_{k=0}^{n} |p_k(t)|.
        \end{aligned}
    \]
    Taking the supremum over all $t \in [0,1]$ on the left side, we obtain the upper bound:
    \[
        \|L_n\| \le \sup_{0\le t\le 1} \sum_{k=0}^{n} |p_k(t)|.
    \]
    To prove this bound is tight, let $\tau \in [0,1]$ be the specific point where the sum attains its global maximum:
    \[
        \sum_{k=0}^{n} |p_k(\tau)| = \sup_{0\le t\le 1} \sum_{k=0}^{n} |p_k(t)| \quad \text{[cite: 76]}.
    \]
    Construct a piecewise linear continuous function $x_0 \in C[0,1]$ such that it precisely captures the signs of the basis polynomials at the nodes:
    \[
        x_0(t_k) = \operatorname{sgn} p_k(\tau) \quad \text{for} \quad 0 \le k \le n \quad \text{[cite: 81]},
    \]
    and ensure that $\|x_0\|_\infty = 1$. Evaluating the interpolation polynomial at $t = \tau$ gives:
    \[
        \begin{aligned}
            L_n x_0(\tau) & = \sum_{k=0}^{n} x_0(t_k) p_k(\tau)                      \\
                          & = \sum_{k=0}^{n} \operatorname{sgn}(p_k(\tau)) p_k(\tau) \\
                          & = \sum_{k=0}^{n} |p_k(\tau)|.
        \end{aligned}
    \]
    This exact match proves that the operator norm achieves the bound:
    \[
        \|L_{n}\| = \sup_{0\le t\le 1}\sum_{k=0}^{n}|p_{k}(t)| \quad \text{[cite: 73]}.
    \]

    \textbf{Part (ii): Application of the Resonance Theorem} \medskip

    We are given the classical lower bound for the Lebesgue constants, which grow at least logarithmically:
    \[
        \sup_{0\le t\le 1}\sum_{k=0}^{n}|p_{k}^{n}(t)| \ge C \log n \quad \text{[cite: 84]}.
    \]
    This strictly implies that the sequence of operator norms diverges to infinity:
    \[
        \sup_{n \ge 0} \|L_n\| = \infty.
    \]

    Now, consider the space $C[0,1]$ equipped with the supremum norm, which is a complete Banach space. The family of bounded linear operators $\{L_n\}$ maps $C[0,1]$ into itself.

    Because the operator family is \textbf{not uniformly bounded}, we apply the contrapositive of the Uniform Boundedness Principle (Banach-Steinhaus Theorem). It dictates that the family cannot be pointwise bounded on the entire Banach space.

    Consequently, there must exist at least one ``resonant" continuous function $x \in C[0,1]$ such that the sequence of its interpolations is unbounded[cite: 86]:
    \[
        \sup_{n\ge 0} \|L_n x\| = \infty \quad \forall n\in\mathbb{N} \quad \text{[cite: 87, 88]}.
    \]
    \textit{Remark:} This rigorously demonstrates that Lagrange interpolating polynomials do not universally converge uniformly to the original function, completely shattering the intuition that ``more nodes always yield better approximations"[cite: 89].
\end{proof}

\subsection*{Reflexive space and best approximation}

\textbf{Intuition and Terminology for Beginners:}
Why does adding the ``Reflexive" condition magically guarantee the existence of a best approximation?
% \begin{itemize}
%     \item \textbf{The Missing Point Problem:} In infinite-dimensional spaces, a sequence of points can wander around a bounded region without ever converging to a solid point (the closed unit ball is not compact). So, a sequence moving closer and closer to the subspace \(E\) might ``vanish into thin air."
%     \item \textbf{Reflexivity to the Rescue:} In a reflexive Banach space, every bounded sequence has a ``weakly convergent" subsequence. It means that even if it doesn't converge in the standard distance (norm), it still settles down to a specific limit point \(y_0\) from the perspective of all linear functionals.
%     \item \textbf{Mazur's Lemma:} If a subspace \(E\) is closed in the normal sense, it is also closed in the weak sense. So the weak limit \(y_0\) cannot escape \(E\).
%     \item \textbf{Weak Lower Semicontinuity:} When a sequence weakly converges to \(y_0\), the norm (length) of the limit point is less than or equal to the eventual lengths of the sequence elements. This ensures that the distance at \(y_0\) exactly hits the minimum.
% \end{itemize}

\begin{theorem}[Kakutani's Theorem on Reflexivity]
    Let $X$ be a Banach space, and let $B_X = \{x \in X : \|x\| \le 1\}$ be its closed unit ball.
    The space $X$ is reflexive if and only if $B_X$ is compact in the weak topology (i.e., weakly compact).
\end{theorem}

\begin{theorem}[Eberlein-\v{S}mulian Theorem - General Topological Version]
    Let $X$ be any Banach space (not necessarily reflexive), and let $A \subset X$ be a subset.
    The following two statements are strictly equivalent:
    \begin{enumerate}
        \item $A$ is weakly compact.
        \item $A$ is weakly sequentially compact (meaning every sequence in $A$ has a subsequence that weakly converges to an element in $A$).
    \end{enumerate}
\end{theorem}

\begin{corollary}[Weak Convergence of Bounded Sequences in Reflexive Spaces]
    Let $X$ be a reflexive Banach space. Every bounded sequence in $X$ contains a weakly convergent subsequence.
\end{corollary}

% \begin{proof}
%     (The proof follows via the separability reduction and Cantor's diagonalization argument as detailed in the previous section. By extracting a countable dense subset of the dual space of the separable span, one constructs a diagonal subsequence that converges for all bounded linear functionals, thus converging weakly.)
% \end{proof}

\begin{lemma}[Mazur's Theorem for Closed Convex Sets]
    Let \(X\) be a Banach space and \(E \subset X\) be a convex set.
    If \(E\) is strongly closed (closed in the norm topology), then \(E\) is also weakly closed.
    (Consequently, if a sequence \(\{z_n\} \subset E\) weakly converges to \(y_0\), then \(y_0 \in E\)).
\end{lemma}

% \begin{proof}
%     Let \(\{x_n\}_{n=1}^\infty\) be a bounded sequence in \(X\), so \(\|x_n\| \le M\) for all \(n\).

%     \textbf{Step 1: Reduce to a separable space.} \\
%     Let \(Y = \overline{\mathrm{span}}\{x_n\}_{n=1}^\infty\). Since \(Y\) is the closure of the span of a countable set, \(Y\) is a separable Banach space. Furthermore, a closed subspace of a reflexive space is also reflexive, so \(Y\) is both separable and reflexive.
%     A well-known theorem in functional analysis states that if a Banach space is reflexive and separable, its dual space \(Y^*\) is also separable. Let \(\{f_k\}_{k=1}^\infty\) be a countable dense subset of \(Y^*\).

%     \textbf{Step 2: Cantor's Diagonalization.} \\
%     Consider the first functional \(f_1\). The sequence of scalars \(\{f_1(x_n)\}_{n=1}^\infty\) is bounded because \(|f_1(x_n)| \le \|f_1\| \|x_n\| \le M \|f_1\|\). By the Bolzano-Weierstrass theorem, it has a convergent subsequence. Let's denote the indices of this subsequence as \(\mathbb{N}_1\).

%     Now consider \(f_2\). We look \textit{only} at the indices in \(\mathbb{N}_1\). The sequence \(\{f_2(x_n)\}_{n \in \mathbb{N}_1}\) is bounded, so it has a convergent subsequence. Let these new indices be \(\mathbb{N}_2 \subset \mathbb{N}_1\).

%     We continue this process infinitely. For each \(k\), we obtain a nested sequence of indices \(\mathbb{N}_1 \supset \mathbb{N}_2 \supset \mathbb{N}_3 \supset \dots\) such that \(\{f_k(x_n)\}_{n \in \mathbb{N}_k}\) converges.

%     We now form the ``diagonal" sequence of indices by choosing \(n_1\) from \(\mathbb{N}_1\), \(n_2\) from \(\mathbb{N}_2\) (with \(n_2 > n_1\)), and so on. For this diagonal subsequence \(\{x_{n_m}\}\), the scalar sequence \(f_k(x_{n_m})\) converges as \(m \to \infty\) for \textbf{every} \(k = 1, 2, 3, \dots\).

%     \textbf{Step 3: Extend convergence to all of \(Y^*\) (The \(\epsilon/3\) argument).} \\
%     Let \(f \in Y^*\) be arbitrary. Let \(\epsilon > 0\). Because \(\{f_k\}\) is dense in \(Y^*\), we can choose an \(f_K\) such that \(\|f - f_K\| < \frac{\epsilon}{3M}\).
%     Since \(f_K(x_{n_m})\) converges, it is a Cauchy sequence. Thus, there exists \(N\) such that for all \(i, j > N\):
%     \[
%         |f_K(x_{n_i}) - f_K(x_{n_j})| < \frac{\epsilon}{3}.
%     \]
%     We now estimate the difference for our arbitrary \(f\):
%     \begin{align*}
%         |f(x_{n_i}) - f(x_{n_j})| & \le |f(x_{n_i}) - f_K(x_{n_i})| + |f_K(x_{n_i}) - f_K(x_{n_j})| + |f_K(x_{n_j}) - f(x_{n_j})|              \\
%                                   & \le \|f - f_K\| \|x_{n_i}\| + \frac{\epsilon}{3} + \|f_K - f\| \|x_{n_j}\|                                 \\
%                                   & < \left(\frac{\epsilon}{3M}\right) M + \frac{\epsilon}{3} + \left(\frac{\epsilon}{3M}\right) M = \epsilon.
%     \end{align*}
%     This shows that for \textit{any} \(f \in Y^*\), the sequence \(\{f(x_{n_m})\}\) is a Cauchy sequence of scalars, hence it converges.
%     This means \(\{x_{n_m}\}\) is a \textit{weakly Cauchy} sequence in \(Y\).

%     \textbf{Step 4: Weak completeness.} \\
%     Since \(Y\) is reflexive, it is weakly complete. This guarantees that the weakly Cauchy sequence \(\{x_{n_m}\}\) actually weakly converges to some element \(y_0 \in Y\).

%     \textbf{Step 5: Lift back to \(X^*\).} \\
%     Finally, for any functional \(g \in X^*\), its restriction to \(Y\) is a bounded linear functional on \(Y\) (i.e., \(g|_Y \in Y^*\)). Therefore, \(g(x_{n_m}) \to g(y_0)\) for all \(g \in X^*\).
%     This completes the proof that \(\{x_{n_m}\}\) weakly converges to \(y_0\) in \(X\).
% \end{proof}

\begin{application}[Existence of Best Approximation in Reflexive Spaces]
    Let \(X\) be a reflexive Banach space, \(E\) a closed subspace of \(X\), and \(x_0 \in X \setminus E\). Let \(d = \mathrm{dist}(x_0, E) > 0\).

    Then, there exists a best approximation \(y_0 \in E\) such that
    \[
        \|x_0 - y_0\| = d.
    \]
    (Furthermore, combining this with the Hahn-Banach distance formula, there exists \(f \in X^*\) with \(\|f\|=1\), \(f|_E = 0\), such that \(f(x_0 - y_0) = \|x_0 - y_0\| = d\)).
\end{application}

\begin{proof}
    By the definition of the infimum \(d = \inf_{z \in E} \|x_0 - z\|\), there exists a minimizing sequence \(\{z_n\}_{n=1}^\infty \subset E\) such that
    \[
        \lim_{n \to \infty} \|x_0 - z_n\| = d.
    \]
    Since the sequence of real numbers \(\|x_0 - z_n\|\) is convergent, it is bounded. By the triangle inequality, \(\|z_n\| \le \|z_n - x_0\| + \|x_0\|\), which implies that the sequence \(\{z_n\}\) is bounded in the Banach space \(X\).

    Because \(X\) is a reflexive space, the Eberlein-\v{S}mulian theorem states that closed bounded sets are weakly sequentially compact. Therefore, the bounded sequence \(\{z_n\}\) must contain a weakly convergent subsequence, denoted as \(\{z_{n_k}\}\). Suppose it weakly converges to a limit \(y_0 \in X\):
    \[
        z_{n_k} \rightharpoonup y_0 \quad \text{as } k \to \infty.
    \]

    Since \(E\) is a closed linear subspace, it is convex and strongly closed. By Mazur's Theorem, a strongly closed and convex set is also weakly closed. Because \(\{z_{n_k}\} \subset E\) and it weakly converges to \(y_0\), we must have
    \[
        y_0 \in E.
    \]

    Now, we consider the sequence \(\{x_0 - z_{n_k}\}\). It is easy to see that it weakly converges to \(x_0 - y_0\). A fundamental property of the norm in Banach spaces is that it is weakly lower semicontinuous (w.l.s.c.). This implies that the norm of the weak limit cannot exceed the limit inferior of the norms of the sequence:
    \[
        \|x_0 - y_0\| \le \liminf_{k \to \infty} \|x_0 - z_{n_k}\|.
    \]

    Since we already know that \(\lim_{k \to \infty} \|x_0 - z_{n_k}\| = d\), we obtain
    \[
        \|x_0 - y_0\| \le d.
    \]
    However, since \(y_0 \in E\), by the definition of the distance \(d\), we naturally have \(\|x_0 - y_0\| \ge d\).

    Combining both inequalities yields
    \[
        \|x_0 - y_0\| = d,
    \]
    which proves that \(y_0 \in E\) is indeed the best approximation of \(x_0\) in \(E\).
\end{proof}











\end{document}